\section*{รายละเอียดผลงาน}
\itemheading{1}{ปัญหา/สาเหตุ/ที่มา เหตุผล และความจำเป็น}
การลานิสิตเกี่ยวข้องกับข้อมูลหลายด้าน ทั้งตัวนิสิต ที่ปรึกษา รายวิชา กลุ่มเรียน ผู้สอน ตารางเรียน และผู้รับผิดชอบการพิจารณา หากข้อมูลเหล่านี้จัดการแยกจากกัน ผู้ใช้ต้องค้นหาความสัมพันธ์และตรวจสอบข้อมูลหลายครั้ง อีกทั้งการติดตามใบลาและสรุปผลกระทบต่อการเรียนต้องใช้ความเข้าใจขั้นตอนร่วมกัน ปัญหาที่โครงการมุ่งแก้จึงเป็นทั้งการเชื่อมข้อมูลและการทำให้กฎการทำงานสามารถอธิบาย ตรวจสอบ และถ่ายทอดได้

โครงการพัฒนาระบบสารสนเทศเป็นศูนย์กลางของคำขอลา โดยนำความรู้เรื่องแบบฟอร์ม ข้อมูลประกอบ ลำดับผู้รับผิดชอบ และการนับวันมาจัดรูปแบบเป็นข้อมูลและกฎที่ใช้ซ้ำได้ แยกระบบดูแลข้อมูลสนับสนุนเพื่อให้ความสัมพันธ์รายวิชา ผู้สอน และที่ปรึกษามีจุดดูแลที่ชัดเจน พร้อมจัดเก็บบริบทใบลาในขณะยื่นเพื่อประกอบการอ่านประวัติ

ในด้านการบริหาร การมีใบลาอยู่ในฐานข้อมูลเพียงอย่างเดียวยังไม่เพียงพอ ต้องมีนิยามจำนวนใบลา วันลา และผลกระทบรายวิชาที่ตรงกัน จึงออกแบบคลังข้อมูลและรายงานตามบทบาท เพื่อเปลี่ยนรายการงานประจำให้เป็นข้อมูลสำหรับติดตามนิสิตและภาพรวมหน่วยงาน โดยควบคุมสิทธิ์และการนับซ้ำ องค์ความรู้ที่ได้สามารถถ่ายทอดผ่านคู่มือ โครงสร้างข้อมูล และเกณฑ์ตรวจสอบ แทนการพึ่งความเข้าใจของผู้ปฏิบัติงานรายบุคคล

\itemheading{2}{วัตถุประสงค์ / เป้าหมาย}
\begin{enumerate}[leftmargin=8mm]
  \item จัดให้มีช่องทางกลางสำหรับยื่น ตรวจสอบ และติดตามใบลาตามหน้าที่ของผู้ใช้
  \item จัดการข้อมูลสนับสนุนและบริบทใบลาให้เชื่อมโยงกัน และตรวจสอบย้อนหลังได้
  \item สังเคราะห์กฎการดำเนินงานและนิยามตัวชี้วัดเป็นองค์ความรู้ที่ใช้และถ่ายทอดร่วมกัน
  \item สนับสนุนรายงานสำหรับนิสิต ที่ปรึกษา ผู้สอน หัวหน้าภาค ผู้บริหาร และเจ้าหน้าที่
  \item วางกระบวนการทวนสอบและปรับปรุงอย่างต่อเนื่อง โดยวัดความถูกต้องของข้อมูล สิทธิ์ และประโยชน์ต่อผู้ใช้ก่อนขยายผล
\end{enumerate}

\itemheading{3}{นวัตกรรมที่นำมาใช้แก้ปัญหาหรือดำเนินการ}
แนวทางของผลงานคือเชื่อมสามส่วนเข้าด้วยกัน ได้แก่ งานใบลา ข้อมูลสนับสนุน และคลังข้อมูล ระบบเว็บแยกการทำงานตามบทบาท บริการฝั่งเซิร์ฟเวอร์ตรวจข้อมูลและสิทธิ์ก่อนทำรายการ กฎการนับวันและลำดับพิจารณาถูกกำหนดในบริการ พร้อมจัดเก็บบริบทข้อมูลขณะยื่นเป็น snapshot เพื่อช่วยอธิบายประวัติของคำขอ

ระบบข้อมูลสนับสนุนดูแลรายวิชา กลุ่มเรียน ตารางเรียน ผู้สอนร่วม และความสัมพันธ์ที่ปรึกษากับนิสิต ส่วนคลังข้อมูลใช้ ETL จัดรูปข้อมูลและแยกข้อเท็จจริงระดับใบลาจากระดับใบลากับกลุ่มเรียน ทำให้วิเคราะห์ได้หลายมุมโดยลดความเสี่ยงการนับซ้ำ รายงานจึงใช้โครงสร้างกลางร่วมกัน แต่จำกัดขอบเขตและรายละเอียดตามหน้าที่

คุณค่าด้านการจัดการความรู้คือการทำให้ความรู้เรื่องขั้นตอน ความสัมพันธ์ข้อมูล และนิยามรายงานอยู่ในรูปที่ตรวจสอบและใช้ซ้ำได้ ทั้งนี้การมีกฎในโค้ดไม่ได้แทนการรับรองระเบียบหน่วยงาน และหน้ารายงานที่ตรวจยังใช้ข้อมูลจำลอง จึงต้องทวนสอบก่อนนำไปอ้างผลการใช้งานจริง
